\subsection{CCNet：Wikipedia 风格的质量代理}

在 Common Crawl 之上，各种后续数据集做了不同程度的清洗。CCNet 的核心思路是：先去重，再做语言识别，最后用 Wikipedia 训练出的 n-gram 模型做质量过滤。CCNet 对低资源语言特别有价值——它的目标不仅是清洗英文网页，还包括为乌尔都语等低资源语言提供高质量预训练数据。

CCNet 的三步流水线：
\begin{enumerate}
    \item \textbf{段落级去重}：对文本做轻量规范化后，移除重复段落。
    \item \textbf{语言识别}：使用 fastText 分类器识别语言，只保留目标语言。
    \item \textbf{质量过滤}：用 Wikipedia 训练的 KenLM 5-gram 模型打分，保留”看起来像 Wikipedia”的文档。
\end{enumerate}

\begin{importantbox}{CCNet 体现了什么}
它体现的是一种非常普遍的数据工程方法：你先定义一个”像好数据的样子”，再把海量候选文本往这个方向过滤。
\end{importantbox}

\subsubsection{Worked Example：基于 Perplexity 的质量过滤}

CCNet 使用的 Wikipedia perplexity 过滤是理解数据质量分类器的最好入口。其训练过程如下：

\begin{enumerate}
    \item \textbf{正例}：取 Wikipedia 的全部文本作为”高质量”的代理分布。
    \item \textbf{训练语言模型}：在 Wikipedia 文本上训练一个 KenLM 5-gram 语言模型。
    \item \textbf{计算 perplexity}：对 Common Crawl 中的每个文档，用该语言模型计算 perplexity。
    \item \textbf{过滤}：perplexity 越低，说明该文档越”像 Wikipedia”，越可能是高质量文本。按 perplexity 分桶，取最低的若干桶作为训练数据。
\end{enumerate}

直觉上，一段写得通顺、信息密集、结构清晰的文本，在 Wikipedia 训练的语言模型下 perplexity 较低；而一段充满广告、重复短语和语法错误的网页，perplexity 则会很高。

\begin{warningbox}{质量过滤的偏见问题}
以 Wikipedia 作为”高质量”代理分布，隐含着一个重要假设：Wikipedia 的写作风格代表了”好文本”。但这个假设并不总是成立：
\begin{itemize}
    \item \textbf{方言与口语}：非标准英语（如 African American Vernacular English）在 Wikipedia 中几乎不出现，perplexity 过滤会系统性地排除这类文本。
    \item \textbf{特定社区}：Reddit 的技术讨论、StackOverflow 的问答、医学论坛的患者叙述——这些对训练有价值的文本，其风格与 Wikipedia 差异很大。
    \item \textbf{创意写作}：诗歌、小说片段、实验性写作的 perplexity 可能很高，但它们对语言模型的表达能力至关重要。
\end{itemize}
RefinedWeb 的作者正是出于这个原因，选择了基于规则的过滤而非基于分类器的过滤——他们认为 ML 过滤器会引入难以察觉的系统性偏见。
\end{warningbox}

